\documentclass[10pt,a4paper]{amsart}
\usepackage[margin=19mm]{geometry}
\usepackage{amsmath,amssymb,mathtools}
\usepackage[scr=rsfs,cal=euler]{mathalfa}
\usepackage{xfrac}
\usepackage[unicode,hidelinks,bookmarks=false]{hyperref}
\newcommand{\ie}{i.e.\ }
\newcommand{\cf}{cf.\ }
\newcommand{\resp}{respectively\ }
\newcommand{\Subsection}{\subsection}
\providecommand{\nobreakdash}{\nobreak\hbox{-}}
\setlength{\emergencystretch}{2em}
\multlinegap0pt
\usepackage{amssymb}
\usepackage{mathtools}
\usepackage{xcolor}
\newtheorem{lemma}{Lemma}
\newtheorem{proposition}{Proposition}
\newcommand{\N}{\mathbb{N}}
\newcommand{\R}{\mathbb{R}}
\newcommand{\Z}{\mathbb{Z}}
\newcommand{\cS}{\mathcal{S}}
\title{Sharp Bounds for Discrete Cube Skeletons}
\author{Dean Menezes}
\date{Source: arXiv:2607.15502v2 (CC BY 4.0)}
\begin{document}
\maketitle

\noindent\emph{Source and licence.} Sharp Bounds for Discrete Cube Skeletons. Version arXiv:2607.15502v2 (CC BY 4.0), distributed by arXiv under the Creative Commons Attribution 4.0 International licence, http://creativecommons.org/licenses/by/4.0/, as recorded in the arXiv licence metadata for this exact version and shown on the abstract page https://arxiv.org/abs/2607.15502v2. Full licence notice: \texttt{licenses/arxiv-2607.15502-CC-BY-4.0.txt}.\par\smallskip

\noindent\textit{Editorial scope.} Counted unit: the article's proposition prop:lower (source0.tex lines 312--319). The article's complete original proof of it is delivered with it (source0.tex lines 321--410, one \texttt{proof} environment of 90 lines). No mathematics is written, repaired or re-derived: the statement and the proof are the article's own lines, in the article's own notation, and no retained line is substituted or re-flowed.  Retained uncounted as the source-local context the proof needs: lines 254--265.  Nothing was omitted from any retained span, so the package carries no omission record.  No retained line contains a citation, so the wrapper carries no bibliography and no bibliography entry.  No retained line contains a figure environment, an image-inclusion command or drawing code, so no image is delivered and the package declares no asset dependency.  Editorial formatting only, in addition: an AMS \texttt{amsart} preamble that restates the article's own declarations and macros used by the retained lines, namely the environments \texttt{lemma}, \texttt{proposition} on separate counters, together with the standard packages those lines need.  The retained environments print in this document as Lemma 1, Proposition 1 in order of appearance, whereas the article prints the same items with the numbering of its own full document.  The complete native source of the article as distributed in the arXiv e-print for this exact version is bundled unchanged as \texttt{sources/205397/source0.tex}.
\begin{lemma}[Fixed-coordinate bound]\label{lem:fixed}
Fix $1\leq\ell\leq n$.  Let $A\subset\R^\ell$ and $T\subset\R^n$ be finite.
Suppose that for every $x\in T$, there is an $r_x>0$ such that
\[
 x_I+r_x\sigma\in A
\]
for every $I\in\binom{[n]}{\ell}$ and $\sigma\in\{\pm1\}^\ell$.  Then
\[
 |A|\geq
 2^{\ell/(2n)}|T|^{\ell(2n-1)/(2n^2)}.
\]
\end{lemma}

\begin{proposition}\label{prop:lower}
For each $0\leq k<n$, there is a constant $c_{n,k}>0$ such that
\[
 |B|\geq c_{n,k}|S|^\beta
\]
whenever finite sets $B,S\subset\Z^n$ and radii $r_x\in\N$ satisfy
$\cS_k(x,r_x)\subset B$ for every $x\in S$.
\end{proposition}

\begin{proof}
Choose $a\in(0,1/2)$ so small that, with $\rho=(2a)^n$,
\[
 2\cdot3^n\rho^{1-\beta}\leq1.
\]
Then $\rho<1$.  Set
\[
 d=\binom{n}{\ell}^{-1}2^{\ell/(2n)-\gamma},
 \qquad
 c_{n,k}=\min\{a^{n\beta},da^k\}.
\]
We prove the claim by strong induction on $N=|S|$.  The case $N=0$ is
trivial.
Put $L=aN^{1/n}$.  If $L<1$, then each skeleton is nonempty and
\[
 |B|\geq1>L^{n\beta}=a^{n\beta}N^\beta\geq c_{n,k}N^\beta.
\]
Assume now that $L\geq1$ and that the claim holds for fewer than $N$ centers.
Call $r_x$ large when $r_x\geq L$ and small otherwise.  At least one class has
at least $N/2$ centers.

Suppose first that the large class $S_{\mathrm L}$ has at least $N/2$ centers.
For $I\in\binom{[n]}{\ell}$, let
\[
 A_I=\{x_I+r_x\sigma:x\in S_{\mathrm L},\
       \ \sigma\in\{\pm1\}^\ell\},
 \qquad
 A=\bigcup_{|I|=\ell}A_I.
\]
Lemma~\ref{lem:fixed} and averaging give an $I$ such that
\[
 |A_I|\geq dN^\gamma.
\]
For each $u\in A_I$, choose $x\in S_{\mathrm L}$ and
$\sigma\in\{\pm1\}^\ell$ with $u=x_I+r_x\sigma$.  With
$J=[n]\setminus I$, the skeleton centered at $x$ contains
\[
 Q_u=\{y\in\Z^n:y_I=u,\ |y_j-x_j|\leq r_x\text{ for }j\in J\}.
\]
The sets $Q_u$ are disjoint because their $I$-coordinate tuples differ.  Each
has $(2r_x+1)^k\geq L^k$ points.  Therefore
\[
 |B|\geq |A_I|L^k
 \geq da^kN^{\gamma+k/n}
 \geq c_{n,k}N^\beta.
\]

Suppose instead that the small class $S_{\mathrm s}$ has at least $N/2$
centers.  Let $h=\lceil L\rceil$ and partition $\Z^n$ into the cells
\[
 P_m=hm+\{0,1,\ldots,h-1\}^n,
 \qquad m\in\Z^n.
\]
Set
\[
 S_m=S_{\mathrm s}\cap P_m,
 \qquad
 N_m=|S_m|,
 \qquad
 Y_m=\bigcup_{x\in S_m}\cS_k(x,r_x).
\]
Since $L\geq1$, we have $h\leq2L$, and hence
\[
 N_m\leq h^n\leq(2L)^n=\rho N<N.
\]
The induction hypothesis gives $|Y_m|\geq c_{n,k}N_m^\beta$ whenever
$N_m>0$.

Each point lies in at most $3^n$ sets $Y_m$.  Indeed, if $y\in Y_m$, then
$|y_i-x_i|<h$ for some $x\in P_m$ and every $i$; for fixed $y_i$, at most
three cells are possible in that coordinate.  Since $\bigcup_m Y_m\subset B$,
\[
 3^n|B|\geq\sum_m |Y_m|
 \geq c_{n,k}\sum_m N_m^\beta.
\]
Since $N_m\leq\rho N$, $\beta<1$, and
$\sum_m N_m=|S_{\mathrm s}|\geq N/2$,
\[
 \sum_m N_m^\beta
 \geq(\rho N)^{\beta-1}\sum_m N_m
 \geq\frac{N^\beta}{2\rho^{1-\beta}}.
\]
Therefore
\[
 |B|\geq
 \frac{c_{n,k}}{2\cdot3^n\rho^{1-\beta}}N^\beta
 \geq c_{n,k}N^\beta.
\]
This completes the induction.
\end{proof}

\end{document}
